% 
\section{The ScaleMaster Dataset}
\label{sec:dataset}

To address the inadequacy of existing benchmarks for evaluating robustness for scale consistency in monocular visual \ac{SLAM}, we introduce and publicly release a new challenging dataset, \textbf{ScaleMaster Dataset}.

% 
\subsection{Hardware}
\label{sec:dataset_hw}

For data acquisition, we built a custom handheld rig (see \figref{fig:system_overview}) equipped with an iPhone 14 Pro, a Livox HAP LiDAR sensor, and an Orbbec Gemini 335L camera. This setup ensures precise temporal synchronization between visual frames and LiDAR scans, leading to improved data quality and reliability. The iPhone provides ARKit odometry, while the LiDAR sensor captures dense 3D geometry. This hardware configuration enables the collection of long, large-scale trajectories with consistent temporal alignment.

% 
\subsection{Baseline Pose Generation}
\label{sec:dataset_baseline}

The baseline camera trajectories were obtained using Apple’s ARKit framework, which is known to provide centimeter-level accuracy in large indoor spaces. Since our primary goal is to highlight the fundamental scale failures, ARKit trajectories that may potentially exhibit long-term drift are sufficient, where residual centimeter-scale errors are tolerable. To build reference maps, we projected high-resolution LiDAR point clouds onto the ARKit trajectories. As shown in \figref{fig:system_overview}, this process yields dense maps that are qualitatively well aligned and preserve fine structural details, making them a reliable reference for analyzing scale consistency and geometric distortion.

% 
\subsection{Dataset Summary and Comparison}
\label{sec:dataset_summary}

The ScaleMaster dataset comprises 25 sequences that specifically address scale-related challenges rarely captured in previous benchmarks. Of these, 18 sequences, acquired solely with an iPhone 14 Pro, provide trajectory-only ground truth, while the remaining 7 sequences, collected using the rig shown in \figref{fig:system_overview}, also include dense LiDAR-based 3D maps. The dataset covers multi-floor motion, long loops with repetitive structures, large open spaces, and low-texture conditions.

\begin{itemize}
  \item \textbf{\tabref{dataset_comparison}} compares ScaleMaster against standard benchmarks, highlighting its unique support for scale evaluation in complex indoor environments.

  \item \textbf{\tabref{dataset_summary}} presents representative sequences with their scale, duration, and environmental tags, illustrating the dataset’s diversity—from nearly 900 m building-scale loops to vertical stair traversals.

\end{itemize}

Each sequence is named by its location followed by an index number (e.g., Library\_01, Parking\_02). The environment description and trajectory characteristics for each sequence are summarized as tags in \tabref{dataset_summary}, covering attributes such as trajectory type, vertical motion, repetitive views, and low-texture conditions. These characteristics establish ScaleMaster as a challenging benchmark for testing SLAM robustness to intra- and inter-session scale inconsistency.



